\documentclass[11pt,a4paper]{article}
\usepackage[margin=2.1cm]{geometry}
\usepackage[T1]{fontenc}
\usepackage{lmodern}
\usepackage{amsmath,amssymb}
\usepackage[hidelinks]{hyperref}
\pagestyle{empty}
\setlength{\parskip}{0.5em}
\setlength{\parindent}{0pt}

\begin{document}

\noindent Carlos Echeverr\'{\i}a-Arrondo\\
Instituto de Dise\~no para la Fabricaci\'on y Producci\'on Automatizada,
MALTA Consolider Team\\
Universitat Polit\`ecnica de Val\`encia, 46022 Val\`encia, Spain\\
\href{mailto:cecharr@idf.upv.es}{cecharr@idf.upv.es}

\vspace{1em}
\noindent Val\`encia, \today

\vspace{1em}
\noindent The Editors\\
\emph{Classical and Quantum Gravity}

\vspace{1.5em}
\noindent Dear Editors,

I am pleased to submit the manuscript ``Boost-range scaling of the
fluctuations of the Benincasa--Dowker action and its consequences for
Everpresent $\Lambda$'' for consideration as a Paper in \emph{Classical and
Quantum Gravity}.

The manuscript answers a question raised in this journal. Moradi, Yazdi and
Zilh\~ao (\emph{Class. Quantum Grav.} 42 (2025) 045017) developed a formalism
for the variance of the Benincasa--Dowker action under Poisson sprinkling and
asked whether these fluctuations could be the origin of Everpresent $\Lambda$,
the causal-set inspired proposal of a cosmological constant that fluctuates
as $V^{-1/2}$. They did not evaluate the variance. We evaluate it,
analytically and with sprinklings in four dimensions, and find that it does
not have the Poisson scaling. The dominant term is the Poisson noise in the
number of nearly null intervals, which grows with the range of Lorentz boosts
that fits inside the region and scales exactly as $(\rho V)^{2-2/d}$; the
remaining terms are a constant fraction of it in every region we simulate
(causal diamonds, past light cones and FLRW past cones reaching the big bang,
up to $\rho V=25\,600$). The action fluctuation therefore grows as
$(\rho V)^{3/4}$, so that a cosmological constant sourced by it would scale as
$V^{-1/4}$. The scaling is a consequence of Lorentz invariance together with
the finite size of the region, and the geometry only fixes the amplitude.

We then follow the consequence through to data. Reading the action
fluctuation as a fluctuation of the cosmological term, and under stated
assumptions about the strongly smeared action, we confront stochastic
$\Lambda(t)$ histories with Pantheon+SH0ES supernovae, DESI DR2 baryon
acoustic oscillations and Planck distance priors. The measured dynamics is
excluded; within a wider family, the data tolerate only a $\Lambda$ that stays
close to a slowly wandering constant. To our knowledge this is the first
confrontation of Everpresent-type dynamics with the DESI DR2 data. The open
points, in particular the limit of strong smearing, are stated explicitly in
the Discussion.

The paper should interest readers of CQG working on causal sets, on discrete
approaches to quantum gravity and on the phenomenology of the cosmological
constant. Every pass/fail criterion of the numerical work was fixed before
the corresponding runs, and the criteria that failed are reported where they
arise. The code, its automated tests and the saved results are public
(GitHub) and archived on Zenodo (DOI 10.5281/zenodo.23413234); the data
statement gives the details.

The manuscript is original, has not been published and is not under
consideration elsewhere. A preprint will be posted on arXiv (the
endorsement for gr-qc is in progress). I am the sole author and have no competing interests. As stated
in the acknowledgements, the numerical code, the analysis and the drafting
were carried out with the assistance of an AI model (Claude, Anthropic); I
designed the project, reviewed all results and take full responsibility for
the content.

Thank you for considering this manuscript.

\vspace{1.5em}
\noindent Yours sincerely,

\vspace{2.5em}
\noindent Carlos Echeverr\'{\i}a-Arrondo

\end{document}
